\section{Protocol-Constrained Contextual Bandit}
\label{sec:method}

\subsection{Legal DB-LBT Actions}

After queue activation, node $k$ draws
$b_k\sim\mathcal U[0,b_{\mathrm{init}}]$. Once an attempt completes, DB-LBT
sets the recovery counter
\begin{equation}
b_k^{+}=
\begin{cases}
\alpha+i_k, & r_k\bmod\kappa<\beta,\\
\mathcal U[0,m], & \text{otherwise},
\end{cases}
\label{eq:dblbt}
\end{equation}
where $r_k$ is the retransmission count. The deterministic branch turns
interruptions into schedule spacing; the short random branch separates
coincident schedules. We retain $\alpha=11$ and expose only the four recovery
parameters in Table~\ref{tab:mapping}.

\begin{table}[!t]
\caption{Contextual-Bandit Actions and Protocol Meaning}
\label{tab:mapping}
\centering
\footnotesize
\begin{tabularx}{\columnwidth}{@{}c p{15mm}Y@{}}
\toprule
Variable & Values & DB-LBT role \\
\midrule
$\kappa$ & $5,7$ & Recovery-cycle length indexed by retransmissions. \\
$\beta$ & $2,3$ & Early states using interruption-derived spacing. \\
$m$ & $4,6,10$ & Maximum random backoff after repeated failure. \\
$b_{\rm init}$ & $15,31$ & Initial range when a queue becomes active. \\
\midrule
$\alpha$ & fixed 11 & Common idle-space offset. \\
\bottomrule
\end{tabularx}
\end{table}

The Cartesian product contains 24 arms and always satisfies
$\beta<\kappa$. It is centered on the published TMC tuple
$(7,3,6,15)$: $b_{\rm init}=15$ matches the common minimum contention
window, 31 gives a more conservative activation range, and each $m$ is below
$\alpha$. An action supplies only the next counter. CCA thresholds,
AIFS/defer timing, NR-U priority class and MCOT, synchronization, MCS, and
ACK/HARQ processing stay fixed by the technology and scenario.

Each parameter addresses a distinct recovery moment. The activation range
$b_{\rm init}$ controls how aggressively a newly nonempty queue enters an
existing schedule. Parameters $\kappa$ and $\beta$ determine how often early
retransmission states return to interruption-derived spacing, while $m$
controls the random escape after repeated failure. A larger value is not
uniformly better: conservative activation can avoid a burst collision but
adds delay when few nodes are active, and extra random recovery can separate
colliding schedules while weakening deterministic reuse. This tradeoff is the
reason to select complete tuples rather than tune one scalar.

\subsection{Local LinUCB Controller}

Each node retains at most 64 of its own attempts. Its 11-dimensional context
$x_k$ contains success/failure ratios, CCA busy fraction, countdown
interruptions, delay EWMA and P95, queue occupancy, arrivals, retransmissions,
effective-data occupancy, and
$\hat n_k=\operatorname{clip}(1+\bar i_k,1,32)$. No peer identity, peer queue,
global fairness value, or cross-technology message enters the vector.

The action space is discrete, the context is low dimensional, and online
samples are scarce. LinUCB fits this setting with one linear reward model per
arm and an explicit exploration bonus. Within a locally separable regime, the
working approximation is
\begin{equation}
\mathbb E[R_{k,t}\mid x_{k,t}=x,q_{k,t}=q,z]
\simeq x^{\mathsf T}\theta_{k,q}^{(z)}.
\label{eq:linear-reward}
\end{equation}
The model is linear in context within one arm; it is not linear in the four
DB-LBT parameters. Keeping a separate $\theta_{k,q}^{(z)}$ for every tuple
allows parameter interactions to differ across arms. It also avoids imposing a
monotone ordering on $\kappa$, $\beta$, $m$, or $b_{\rm init}$. Thus the
nonlinear CCA and decoding maps in \eqref{eq:cca}--\eqref{eq:success} remain
outside the learner, while the locally observed reward surface is approximated
where decisions are actually made.

For each arm $q$, the controller maintains
$A_q\in\mathbb R^{11\times11}$ and $d_q\in\mathbb R^{11}$:
\begin{align}
\hat\theta_q&=A_q^{-1}d_q,\\
p_q(x_k)&=\hat\theta_q^{\mathsf T}x_k+
c\sqrt{x_k^{\mathsf T}A_q^{-1}x_k},\\
q^*&=\arg\max_q p_q(x_k),
\label{eq:linucb}
\end{align}
with $c=0.5$. After observing reward $R_k$, the selected model receives
$A_q\leftarrow A_q+x_kx_k^{\mathsf T}$ and
$d_q\leftarrow d_q+R_kx_k$.

The reward is local as well. Let $A_k$ be effective-airtime fraction,
$D_{95,k}$ P95 delay, $s_k$ the node's successful-busy share, and $P_{f,k}$ its
failure ratio:
\begin{align}
U_A&=\operatorname{clip}(\hat n_kA_k,0,1),\\
U_D&=1-\operatorname{clip}
 \left(\frac{D_{95,k}}{16\,\mathrm{ms}\,\hat n_k},0,1\right),\\
U_S&=1-\operatorname{clip}
 \left(\frac{|s_k-1/\hat n_k|}{1/\hat n_k},0,1\right),\\
R_k&=(U_A+U_D+U_S)/3-0.25P_{f,k}.
\label{eq:reward}
\end{align}
Every transmitter owns a controller initialized from the same offline model.
After eight observations, the event implementation decides every 32 global
contention rounds; the packet implementation decides every 32 completed local
attempts. A decision affects the next recovery only, never a countdown already
in progress.

\subsection{Initialization and Online Sequence}

Offline initialization evaluates every arm under 11 topology, load,
interference, and sensing conditions using three training seeds. Those samples
populate the 24 linear models before formal evaluation; none of the ten formal
event seeds is used for initialization or fixed-oracle selection. The same
initial model is copied to each transmitter, but subsequent contexts, choices,
rewards, and updates are local.

An online step follows the protocol clock. The node first aggregates its latest
history into $x_k$, selects $q^*$ with \eqref{eq:linucb}, and applies that tuple
only when the current attempt has completed. It then observes the next reward
window and updates the chosen arm. Separating observation, selection,
application, and reward assignment prevents an action from receiving credit
for a countdown that began under an earlier profile.
